\documentclass[11pt,a4paper]{article}
\usepackage[utf8]{inputenc}
\usepackage[T1]{fontenc}
\usepackage{amsmath,amssymb,amsthm,mathtools}
\usepackage[margin=2.5cm]{geometry}
\usepackage{booktabs}
\usepackage{array}
\usepackage{hyperref}
\usepackage[authoryear,round]{natbib}
\usepackage{listings}
\lstset{basicstyle=\small\ttfamily,breaklines=true}
\setlength{\emergencystretch}{2em}

\theoremstyle{plain}
\newtheorem{theorem}{Theorem}[section]
\newtheorem{lemma}[theorem]{Lemma}
\newtheorem{corollary}[theorem]{Corollary}
\newtheorem{proposition}[theorem]{Proposition}
\newtheorem{conjecture}[theorem]{Conjecture}
\newtheorem{conditionaltheorem}[theorem]{Conditional Theorem}
\theoremstyle{definition}
\newtheorem{definition}[theorem]{Definition}
\theoremstyle{remark}
\newtheorem{remark}[theorem]{Remark}

\newcommand{\PCF}{\mathrm{PCF}}
\newcommand{\half}{\tfrac{1}{2}}
\newcommand{\dff}[1]{(#1)!!}
\newcommand{\hyp}[5]{{{}_{#1}F_{#2}}\!\left(\genfrac{}{}{0pt}{}{#3}{#4}\bigg|#5\right)}

\title{Polynomial Continued Fractions: a Proved Logarithmic Ladder,\\
a \texorpdfstring{$4/\pi$}{4/pi} Casoratian Identity, and 482 Irrational Constants}
\author{Papanokechi}
\date{April 2026}

\begin{document}
\maketitle

\begin{abstract}
We study polynomial continued fractions and prove two main theorems
unconditionally.
The first establishes a Logarithmic Ladder: a one-parameter quadratic
family converges to the reciprocal logarithm for every real parameter
greater than one.
The second derives a Casoratian telescoping identity yielding $4/\pi$
and a rapidly convergent series for $\pi/4$ with error
$O(2^{-N}/N^{7/2})$, from the PCF with $\alpha(n) = -n(2n{-}3)$
and $\beta(n) = 3n{+}1$.
We prove that this $4/\pi$ family cannot arise from any Gauss
${}_2F_1$ continued fraction, its series representation being forced
into a ${}_2F_3$ form over $\mathbb{Q}(\sqrt{5})$.
A conditional extension, assuming a conjectural ratio formula
(Conjecture~\ref{conj:ratio}), yields a continuous Gamma-function
identity for all non-negative integers~$m$.
As a computational application, we certify 482~quadratic generalized
continued-fraction constants using Arb ball arithmetic at
$10\,000$-bit precision and depth $N = 5000$, proving irrationality
for each by Euler's criterion.
PSLQ integer-relation searches against 15~standard constants and
588~Gamma-ratio expressions find no match.
\end{abstract}

\noindent
\textbf{2020 Mathematics Subject Classification:}
11J70 (primary), 11Y60, 33C20.

\noindent
\textbf{Keywords:}
continued fractions, irrationality, Casoratian,
Gauss hypergeometric, Pincherle's theorem, PSLQ.

\tableofcontents

%======================================================================
\section{Introduction}\label{sec:intro}
%======================================================================

%--- (a) Motivating context ---

A \emph{polynomial continued fraction} (PCF) with numerator polynomial
$\alpha(n)$ and denominator polynomial~$\beta(n)$ is
\begin{equation}\label{eq:pcf}
  v = \beta(0) + \cfrac{\alpha(1)}{\beta(1) +
  \cfrac{\alpha(2)}{\beta(2) + \cfrac{\alpha(3)}{\beta(3) + \cdots}}}.
\end{equation}
The convergents $C_n = p_n/q_n$ are computed via the forward recurrence
\begin{equation}\label{eq:rec}
  p_n = \beta(n)\,p_{n-1} + \alpha(n)\,p_{n-2}, \qquad
  q_n = \beta(n)\,q_{n-1} + \alpha(n)\,q_{n-2},
\end{equation}
with $p_{-1} = 1$, $p_0 = \beta(0)$, $q_{-1} = 0$, $q_0 = 1$.
Classical work of \citet{wall1948} and \citet{lorentzen2008}
provides the analytic background on continued fractions, convergence criteria,
and the Gauss hypergeometric model relevant here.

The irrationality proof of \citet{apery1979} for $\zeta(3)$ showed
that polynomial continued fractions can carry deep arithmetic content.
More recently, the Ramanujan Machine program
\citep{raayoni2021,elimelech2024} has renewed the search for
low-degree examples with closed-form values, and a companion paper
\citep{papanokechi2026spectral} develops a spectral classification
framework for such families.
In this paper we prove unconditional closed evaluations for two
explicit one-parameter families: one producing a full logarithmic
ladder $1/\ln(k/(k{-}1))$ for every real $k > 1$, and another
yielding $4/\pi$ together with a rapidly convergent series identity
for~$\pi/4$.

%--- (b) Contributions ---

The contributions of this paper are as follows.

\medskip
\noindent\textit{Proved results.}
\begin{enumerate}
\item \textbf{Logarithmic Ladder} (Theorem~\ref{thm:log}).
  The PCF with $\alpha(n) = -kn^2$ and $\beta(n) = (k{+}1)n + k$
  converges to $1/\ln(k/(k{-}1))$ for all real $k > 1$.
\item \textbf{$\boldsymbol{4/\pi}$ Casoratian identity}
  (Theorems~\ref{thm:conv1}, \ref{thm:cas}, and~\ref{thm:4overpi}).
  The PCF with $\alpha(n) = -n(2n{-}3)$, $\beta(n) = 3n{+}1$
  converges to $4/\pi$, and the Casoratian telescope yields the
  series~\eqref{eq:series} with error $O(2^{-N}/N^{7/2})$.
\item \textbf{Gauss CF non-membership} (Proposition~\ref{prop:gauss}).
  The $4/\pi$ family is not a Gauss ${}_2F_1$ continued fraction;
  its series is a ${}_2F_3$ over $\mathbb{Q}(\sqrt{5})$.
\item \textbf{482 irrationality certificates}
  (Proposition~\ref{prop:irrat}).
  All 482 quadratic generalized continued-fraction constants
  $V(A,B,C)$ are proved irrational by Euler's criterion,
  with Arb-certified enclosures at $10\,000$-bit precision
  and PSLQ searches against 15 standard constants finding no match.
\end{enumerate}

\noindent\textit{Conditional results} (assuming
Conjecture~\ref{conj:ratio}).
\begin{enumerate}
\setcounter{enumi}{4}
\item \textbf{Gamma identity}
  (Conditional Theorem~\ref{thm:gamma},
   Corollary~\ref{cor:pi},
   Corollary~\ref{cor:special}).
  For all non-negative integers~$m$, the Pi Family PCF converges to
  $2\,\Gamma(m{+}1)/(\sqrt{\pi}\,\Gamma(m{+}\frac{1}{2}))$,
  extending the base cases $m = 0,1$ (which are proved unconditionally).
\end{enumerate}

\noindent\textit{Numerical/experimental.}
\begin{enumerate}
\setcounter{enumi}{5}
\item \textbf{Cubic Clausen witness}
  (Section~\ref{sec:exp-cq}; numerical/experimental conjecture).
  The cubic J-fraction $\operatorname{eval\_cf}(2,3,2,2/3)$ matches a
  Clausen square ${}_2F_1(u,1{-}u;\frac{3}{2};\frac{27}{64})^2$
  with residual below $10^{-341}$ at 350-digit precision;
  this match is not claimed as a theorem and is recorded as
  Open Question~2.
\end{enumerate}

%--- (c) Proof methods ---

The proofs of the Logarithmic Ladder and the Pi Family base cases
($m = 0,1$) proceed by closed-form induction on the convergent
numerators and denominators, followed by identification of the
limiting ratio via classical series.
The $4/\pi$ evaluation uses a discrete raising operator $L$ that maps
the $m = 0$ solution space to~$m = 1$ (Lemma~\ref{lem:shift}),
combined with Poincar\'e growth analysis.
The Casoratian series identity is derived by a telescoping argument
internal to the $m = 1$ recurrence.
The Gauss non-membership proof reduces to a polynomial system showing
no rational ${}_2F_1$ parameters match.

%--- (d) Computational methodology ---

All numerical results are certified using Arb ball arithmetic via
\texttt{python-flint}~0.8.0 at $10\,000$-bit precision (approximately
$3\,010$ decimal digits) and forward depth $N = 5000$.
The bracketing method of Lemma~\ref{lem:bracket} provides rigorous
enclosures: the main proved families achieve bracket widths below
$10^{-1500}$, yielding over 1500 correct digits.
The 482-constant catalogue was generated by a parametric sweep over
$(A,B,C)$ with $A \in \{1,\ldots,5\}$, $|B| \le 5$, $C \in \{1,\ldots,5\}$,
followed by PSLQ integer-relation searches at 300-digit precision.

The proved statements are separated carefully from the conjectural ones.
Section~\ref{sec:log} establishes the logarithmic ladder for all real $k>1$.
Sections~\ref{sec:pi} and~\ref{sec:closed} prove the base cases $m=0$ and
$m=1$ of the $\pi$-family, and Section~\ref{sec:casoratian} derives the
Casoratian identity for $\pi/4$. Section~\ref{sec:gauss} proves non-membership
in the Gauss continued-fraction class, Section~\ref{sec:vquad} proves
irrationality for the quadratic catalogue, and Section~\ref{sec:exp-cq}
records a cubic Clausen-square witness.
Conjecture~\ref{conj:ratio} proposes the general ratio law in the parameter
$m$, and Section~\ref{sec:gamma} derives the continuous Gamma quotient from
that conjecture together with the proved base cases.
Sections~\ref{sec:euler}--\ref{sec:meta}
place the proved formulas in context, while the appendices record numerical
certification, the search procedure, exact convergents, and reproducibility data.

%======================================================================
\section{The Logarithmic Ladder}\label{sec:log}
%======================================================================

\begin{theorem}[Logarithmic Ladder]\label{thm:log}
For all real $k > 1$, the PCF with
$\alpha(n) = -kn^2$ and $\beta(n) = (k{+}1)n + k$
converges to $1/\ln(k/(k{-}1))$.
\end{theorem}

\begin{proof}
\textbf{Claim.}
\begin{equation}\label{eq:log-pq}
  p_n = (n{+}1)!\, k^{n+1}, \qquad
  q_n = (n{+}1)!\, \sum_{j=0}^{n} \frac{k^{n-j}}{j+1}.
\end{equation}

\textbf{Base case.}
$p_0 = \beta(0) = k = 1!\cdot k^1$ and $q_0 = 1 = 1!\cdot k^0/1$.

\medskip\noindent\textbf{Induction for $p_n$.}
Assume $p_{n-1} = n!\,k^n$ and $p_{n-2} = (n{-}1)!\,k^{n-1}$.
\begin{align*}
  \beta(n)\,p_{n-1} + \alpha(n)\,p_{n-2}
  &= \bigl((k{+}1)n + k\bigr) n!\,k^n - kn^2 (n{-}1)!\,k^{n-1} \\
  &= n!\,k^n\bigl((k{+}1)n + k - n\bigr)
   = n!\,k^n \cdot k(n{+}1) = (n{+}1)!\,k^{n+1}.
\end{align*}

\medskip\noindent\textbf{Induction for $q_n$.}
Write $S_n = \sum_{j=0}^{n} k^{n-j}/(j{+}1)$
and note $S_n = k\,S_{n-1} + 1/(n{+}1)$.
Assume $q_{n-1} = n!\,S_{n-1}$ and $q_{n-2} = (n{-}1)!\,S_{n-2}$.
Using $S_{n-1} = k\,S_{n-2} + 1/n$:
\begin{align*}
\beta(n)\,q_{n-1} + \alpha(n)\,q_{n-2}
&= \bigl((k{+}1)n{+}k\bigr)\,n!\bigl(k\,S_{n-2}+\tfrac{1}{n}\bigr)
   - kn^2\,(n{-}1)!\,S_{n-2} \\
&= S_{n-2}\cdot k^2(n{+}1)\,n!
   + (n{-}1)!\bigl(kn+k+n\bigr).
\end{align*}
Meanwhile,
$(n{+}1)!\,S_n
= k^2(n{+}1)\,n!\,S_{n-2} + k(n{+}1)(n{-}1)! + n!$.
Since $(n{-}1)!(kn{+}k{+}n) = k(n{+}1)(n{-}1)! + n!$,
the two expressions are equal.

\medskip\noindent\textbf{Limit.}
The closed forms~\eqref{eq:log-pq} are polynomial identities in~$k$,
valid for all real~$k$.
The convergent ratio is
$C_n = k/\sum_{j=0}^{n}(1/k)^j/(j{+}1)$.
For $|1/k| < 1$, the series
$\sum_{j=0}^{\infty} x^j/(j{+}1) = -\ln(1{-}x)/x$
converges at $x = 1/k$, giving
$\sum (1/k)^j/(j{+}1) = k\ln(k/(k{-}1))$.
Therefore
$\lim C_n = k/(k\ln(k/(k{-}1))) = 1/\ln(k/(k{-}1))$.
\end{proof}

\begin{remark}
The theorem holds for all real $k > 1$, not just integers.
As an illustration, using Arb ball arithmetic at $120$ decimal digits
and forward depth $N=2000$, we checked $15$ non-integer values in
$[1.5,100]$ and found at least $90$ matching digits with
$1/\ln(k/(k{-}1))$.
Table~\ref{tab:log} lists representative instances.
(This remark is purely illustrative; Theorem~\ref{thm:log} is proved
unconditionally above.)
\end{remark}

\begin{table}[ht]
\centering
\caption{Logarithmic Ladder: first instances}\label{tab:log}
\begin{tabular}{@{}cccc@{}}
\toprule
$k$ & $\alpha(n)$ & $\beta(n)$ & Value \\
\midrule
$2$ & $-2n^2$ & $3n+2$ & $1/\ln 2$ \\
$3$ & $-3n^2$ & $4n+3$ & $1/\ln(3/2)$ \\
$4$ & $-4n^2$ & $5n+4$ & $1/\ln(4/3)$ \\
$5$ & $-5n^2$ & $6n+5$ & $1/\ln(5/4)$ \\
$\varphi$ & $-\varphi\, n^2$ & $(\varphi{+}1)n{+}\varphi$ &
  $1/\ln(\varphi/(\varphi{-}1))$ \\
$k$ & $-kn^2$ & $(k{+}1)n{+}k$ & $1/\ln(k/(k{-}1))$ \\
\bottomrule
\end{tabular}
\end{table}

%======================================================================
\section{The Pi Family}\label{sec:pi}
%======================================================================

\subsection{Base case: \texorpdfstring{$2/\pi$}{2/pi}}\label{sec:pi_base}

\begin{theorem}[Pi Family, $m=0$]\label{thm:pi0}
The PCF with $\alpha(n) = -n(2n{-}1)$ and $\beta(n) = 3n+1$
converges to $2/\pi$.
\end{theorem}

\begin{proof}
\textbf{Claim.}
\begin{equation}\label{eq:pi-pq}
  p_n = (2n{+}1)!!, \qquad
  q_n = (2n{+}1)!! \sum_{j=0}^{n} \frac{j!}{(2j{+}1)!!}.
\end{equation}

\textbf{Base case.}
$p_0 = 1 = (1)!!$ and $q_0 = 1 = 1\cdot(0!/1!!)$.

\medskip\noindent\textbf{Induction for $p_n$.}
Assume $p_{n-1} = (2n{-}1)!!$ and $p_{n-2} = (2n{-}3)!!$.
Using $(2n{-}1)!! = (2n{-}1)(2n{-}3)!!$:
\begin{align*}
  \beta(n)\,p_{n-1} + \alpha(n)\,p_{n-2}
  &= (3n{+}1)(2n{-}1)!! - n(2n{-}1)(2n{-}3)!! \\
  &= (2n{-}3)!!\bigl[(3n{+}1)(2n{-}1) - n(2n{-}1)\bigr] \\
  &= (2n{-}3)!!(2n{-}1)(2n{+}1) = (2n{+}1)!!.
\end{align*}

\medskip\noindent\textbf{Induction for $q_n$.}
Write $T_n = \sum_{j=0}^{n} j!/(2j{+}1)!!$.
Note $T_{n-1} = T_{n-2} + (n{-}1)!/(2n{-}1)!!$.
With inductive hypotheses
$q_{n-1} = (2n{-}1)!!\,T_{n-1}$ and $q_{n-2} = (2n{-}3)!!\,T_{n-2}$:
\begin{align*}
\beta(n)\,q_{n-1} + \alpha(n)\,q_{n-2}
&= (2n{-}1)!!\bigl[(3n{+}1)\,T_{n-1} - n\,T_{n-2}\bigr].
\end{align*}
Substituting $T_{n-1} = T_{n-2} + (n{-}1)!/(2n{-}1)!!$:
\[
  (3n{+}1)\,T_{n-1} - n\,T_{n-2}
  = (2n{+}1)\,T_{n-2} + \frac{(3n{+}1)(n{-}1)!}{(2n{-}1)!!}.
\]
The identity reduces to
$(3n{+}1)(n{-}1)! = (2n{+}1)(n{-}1)! + n!$,
which holds since $n! = n(n{-}1)!$.

\medskip\noindent\textbf{Limit.}
From the closed forms~\eqref{eq:pi-pq}, we have $C_n = p_n/q_n = 1/T_n$.
The identity \citep[\S5.2]{cuyt2008}
\begin{equation}\label{eq:pi-id}
  \frac{\pi}{2}
  = \sum_{j=0}^{\infty} \frac{j!}{(2j{+}1)!!}
  = \sum_{j=0}^{\infty} \frac{2^j(j!)^2}{(2j{+}1)!},
\end{equation}
derivable from $\arcsin(1) = \pi/2$,
gives $\lim T_n = \pi/2$ and $\lim C_n = 2/\pi$.
\end{proof}

\subsection{Parity phenomenon}\label{sec:parity}

\begin{remark}[Parity]\label{thm:parity}
For even $c = 2\ell$, the PCF with $\alpha(n) = -n(2n{-}c)$ and
$\beta(n) = 3n+1$ evaluates to an exact rational, since
$\alpha(\ell) = -\ell(2\ell - 2\ell) = 0$ and the CF terminates
at depth~$\ell{-}1$.
\end{remark}

The first even specialisations are listed in Table~\ref{tab:parity}.

\begin{table}[ht]
\centering
\caption{Pi Family: even parameters (exact rationals)}\label{tab:parity}
\begin{tabular}{@{}rcc@{}}
\toprule
$c$ & Value & Formula \\
\midrule
$2$ & $1$ & $1$ \\
$4$ & $3/2$ & $\binom{2}{1}/2$ \\
$6$ & $15/8$ & $\binom{4}{2}/2^3$ \\
$8$ & $35/16$ & $\binom{6}{3}/2^4$ \\
$10$ & $315/128$ & $\binom{8}{4}/2^7$ \\
$12$ & $693/256$ & $\binom{10}{5}/2^8$ \\
\bottomrule
\end{tabular}
\end{table}

\subsection{Odd parameters and the binomial recurrence}

\begin{conjecture}[Ratio Formula]\label{conj:ratio}
Define $\operatorname{val}(m)$ as the limit of the PCF with
$\alpha(n) = -n(2n{-}(2m{+}1))$ and $\beta(n) = 3n+1$.
Then $\operatorname{val}(m{+}1) = \operatorname{val}(m)\cdot 2(m{+}1)/(2m{+}1)$.
\end{conjecture}

\begin{remark}
Using Arb ball arithmetic at $600$ decimal digits and depth $N=5000$,
the recurrence matches to at least $500$ digits for $m = 0,1,\ldots,10$
and for the half-integer values $m = 0.5,1.5,\ldots,10.5$.
It is consistent with the closed form
$\operatorname{val}(m) = 2^{2m+1}/(\pi\binom{2m}{m})$, since
$\binom{2m+2}{m+1} = \binom{2m}{m}\cdot 2(2m{+}1)/(m{+}1)$ gives
ratio $4\binom{2m}{m}/\binom{2m+2}{m+1} = 2(m{+}1)/(2m{+}1)$.
\end{remark}

\begin{corollary}[Full Pi Family {\normalfont(conditional on Conjecture~\ref{conj:ratio})}]\label{cor:pi}
Assuming Conjecture~\ref{conj:ratio}, for all $m \ge 0$ the PCF
converges to $2^{2m+1}/(\pi\binom{2m}{m})$.
\end{corollary}

\begin{proof}
This is immediate from Theorem~\ref{thm:pi0} and repeated application of
Conjecture~\ref{conj:ratio}. Starting from $\operatorname{val}(0)=2/\pi$,
one obtains
\[
  \operatorname{val}(m)=\frac{2}{\pi}\prod_{j=0}^{m-1}\frac{2(j+1)}{2j+1}
  = \frac{2^{2m+1}}{\pi\binom{2m}{m}},
\]
which is the standard central-binomial product identity.
\end{proof}

Table~\ref{tab:pi} records the first odd specialisations.

\begin{table}[ht]
\centering
\caption{Pi Family: odd parameter $c = 2m{+}1$}\label{tab:pi}
\begin{tabular}{@{}ccccc@{}}
\toprule
$m$ & $c$ & $\alpha(n)$ & $\operatorname{val}(m)\cdot\pi$ &
$\operatorname{val}(m)$ \\
\midrule
$0$ & $1$ & $-n(2n{-}1)$ & $2$ & $2/\pi$ \\
$1$ & $3$ & $-n(2n{-}3)$ & $4$ & $4/\pi$ \\
$2$ & $5$ & $-n(2n{-}5)$ & $16/3$ & $16/(3\pi)$ \\
$3$ & $7$ & $-n(2n{-}7)$ & $32/5$ & $32/(5\pi)$ \\
$4$ & $9$ & $-n(2n{-}9)$ & $256/35$ & $256/(35\pi)$ \\
\bottomrule
\end{tabular}
\end{table}

%======================================================================
\section{Convergent Closed Forms}\label{sec:closed}
%======================================================================

\begin{theorem}[$m=0$ convergent numerators]\label{thm:conv0}
For the Pi Family with $m=0$,
$p_n = (2n{+}1)!! = \dff{2n{-}1}\cdot(2n{+}1)$.
\end{theorem}

\begin{proof}
Base cases: $p_0{=}1$, $p_1{=}3$, $p_2{=}15$.
Induction: assume $p_k = (2k{+}1)!!$ for $k < n$.
The recurrence $p_n = (3n{+}1)p_{n-1} - n(2n{-}1)p_{n-2}$
with $P(n) = 2n{+}1$ reduces to the identity
$(2n{-}1)(2n{-}3)P(n) = (3n{+}1)(2n{-}3)P(n{-}1) - n(2n{-}1)P(n{-}2)$
after dividing by~$\dff{2n{-}3}$.
Expanding both sides yields $0 = 0$.
\end{proof}

\begin{theorem}[$m=1$ convergent numerators]\label{thm:conv1}
For the Pi Family with $m=1$ ($\alpha(n) = -n(2n{-}3)$,
$\beta(n) = 3n{+}1$),
\begin{equation}\label{eq:pn_m1}
  p_n = \dff{2n{-}1}\,(n^2 + 3n + 1).
\end{equation}
The polynomial $h_n = n^2{+}3n{+}1$ \citep{oeis_A028387}
has discriminant~$5$ and
roots $(-3 \pm \sqrt{5})/2$ related to the golden ratio.
\end{theorem}

\begin{proof}
After the substitution $p_n = \dff{2n{-}1}\,u_n$
(using $\dff{2n{-}1} = (2n{-}1)\dff{2n{-}3}$),
the recurrence reduces to
\begin{equation}\label{eq:ured}
  (2n{-}1)\,u_n = (3n{+}1)\,u_{n-1} - n\,u_{n-2}.
\end{equation}
Substituting $u_n = n^2{+}3n{+}1$, $u_{n-1} = n^2{+}n{-}1$,
$u_{n-2} = n^2{-}n{-}1$:
\begin{align*}
  \text{RHS} &= (3n{+}1)(n^2{+}n{-}1) - n(n^2{-}n{-}1) \\
  &= 2n^3 + 5n^2 - n - 1
   = (2n{-}1)(n^2{+}3n{+}1) = \text{LHS.}\qedhere
\end{align*}
\end{proof}

\begin{conjecture}[General $m$]\label{conj:Rnm}
For all $m \ge 0$ and fixed~$n$, the quantity
$R(n,m) = p_n(m)/\dff{2n{-}1}$
is a polynomial in~$m$ of degree $\lceil n/2 \rceil$.
\end{conjecture}

A symbolic computation confirms this for $m = 0, \ldots, 12$ and $n \le 20$
(see the script \texttt{conj43\_check.py} in the repository,
Appendix~\ref{app:repro}).
For $m \ge 2$, $R(n,m)$ is rational with denominators dividing
$\dff{2n{-}1}$.

%======================================================================
\section{The Continuous Gamma Identity}\label{sec:gamma}
%======================================================================

The results of \S\S\ref{sec:log}--\ref{sec:closed} are special
cases of a single continuous identity. Its extension beyond the
base cases $m=0,1$ is conditional on Conjecture~\ref{conj:ratio}.
The base cases $m=0$ and $m=1$ are proved unconditionally in Steps~1--3
below; the extension to general~$m$ in Step~4 assumes Conjecture~\ref{conj:ratio}.

\begin{conditionaltheorem}[Gamma Function Representation; unconditional for $m=0,1$; conditional on Conjecture~\ref{conj:ratio} for $m \geq 2$]\label{thm:gamma}
The formula below is proved unconditionally for $m=0$ and $m=1$.
Assuming Conjecture~\ref{conj:ratio}, for all real $m > -\frac{1}{2}$,
the PCF with
$\alpha(n) = -n(2n{-}2m{-}1)$ and $\beta(n) = 3n+1$
converges to
\[
  \operatorname{val}(m)
  = \frac{2\,\Gamma(m{+}1)}{\sqrt{\pi}\,\Gamma(m{+}\half)}.
\]
\end{conditionaltheorem}

\begin{proof}
We prove the identity for $m = 0$ and $m = 1$ by a purely algebraic
argument, then extend to all~$m$ via the Wallis reduction formula and
Conjecture~\ref{conj:ratio} (supported numerically at $600$ decimal digits and depth $N=5000$).

\medskip
\emph{Step~1: The raising operator.}
Define $L(f)_n = (n{+}2)\,f_n - (n{+}1)^2\,f_{n-1}$.

\begin{lemma}\label{lem:shift}
If\/ $f$ satisfies the $m{=}0$ recurrence
$f_n = (3n{+}1)f_{n-1} - n(2n{-}1)f_{n-2}$,
then $g = L(f)$ satisfies the $m{=}1$ recurrence
$g_n = (3n{+}1)g_{n-1} - n(2n{-}3)g_{n-2}$.
\end{lemma}

\begin{proof}[Proof of Lemma]
Write $g_{n-1} = (n{+}1)f_{n-1} - n^2 f_{n-2}$ and
$g_{n-2} = n\,f_{n-2} - (n{-}1)^2 f_{n-3}$.
Expanding $(3n{+}1)g_{n-1} - n(2n{-}3)g_{n-2}$ and
eliminating $f_{n-3}$ via the $m{=}0$ recurrence at index~$n{-}1$
yields coefficients
\begin{align*}
  \text{coeff.\ of } f_{n-1} &:
    (3n^2{+}4n{+}1) - n(n{-}1) = 2n^2{+}5n{+}1, \\
  \text{coeff.\ of } f_{n-2} &:
    {-}n^2(5n{-}2) + n(n{-}1)(3n{-}2) = {-}n(2n^2{+}3n{-}2),
\end{align*}
matching
$g_n = (n{+}2)[(3n{+}1)f_{n-1} - n(2n{-}1)f_{n-2}]
  - (n{+}1)^2 f_{n-1}$
identically.
\end{proof}

Since $L$ is linear, it maps the solution space~$\mathcal{S}_0$
of the $m{=}0$ recurrence into~$\mathcal{S}_1$.
Checking initial values:
\begin{alignat*}{2}
  L(P^{(0)})_0 &= 2\cdot 1 - 1\cdot 1 = 1 = P_0^{(1)},\;
  &L(P^{(0)})_1 &= 3\cdot 3 - 4\cdot 1 = 5 = P_1^{(1)},\\
  L(Q^{(0)})_0 &= 2\cdot 1 - 1\cdot 0 = 2 = 2\,Q_0^{(1)},\;
  &L(Q^{(0)})_1 &= 3\cdot 4 - 4\cdot 1 = 8 = 2\,Q_1^{(1)}.
\end{alignat*}
By uniqueness:
$P_n^{(1)} = L(P^{(0)})_n$ and $2\,Q_n^{(1)} = L(Q^{(0)})_n$.

\medskip
\emph{Step~2: Poincar\'e growth analysis.}
The recurrence has Poincar\'e characteristic roots $\lambda = 2$
and $\lambda = 1$
(from $\lambda^2 - 3\lambda + 2 = 0$).
Both $P_n^{(m)}$ and~$Q_n^{(m)}$ are dominant, growing
like~$(2n)!!$, so $P_n^{(m)}/Q_n^{(m)} \to S^{(m)}$ at
rate~$O(n!/(2n)!!) = O(2^{-n}/\sqrt{n})$.

\medskip
\emph{Step~3: The limit.}
From Step~1:
\[
  S_n^{(1)}
  = 2\cdot\frac{(n{+}2)P_n^{(0)} - (n{+}1)^2 P_{n-1}^{(0)}}
              {(n{+}2)Q_n^{(0)} - (n{+}1)^2 Q_{n-1}^{(0)}}.
\]
Dividing by $(n{+}2)Q_n^{(0)}$ and using
$(n{+}1)^2/[(n{+}2)(2n{+}1)] \to \frac{1}{2}$,
both correction terms converge to~$\frac{1}{2}$, giving
$S^{(1)} = 2\,S^{(0)}$.

\medskip
\emph{Step~4: General~$m$ and the Gamma identity.}
The base case $S^{(0)} = 2/\pi$ is established by
Theorem~\ref{thm:pi0}.
The Wallis reduction formula gives
$I_{m+1} = \frac{2m{+}1}{2m{+}2}\,I_m$
where $I_m = \int_0^{\pi/2}\sin^{2m}(x)\,dx
= \frac{\sqrt{\pi}\,\Gamma(m{+}\frac{1}{2})}{2\,\Gamma(m{+}1)}$.
Assuming Conjecture~\ref{conj:ratio},
the ratio formula
$\operatorname{val}(m{+}1)/\operatorname{val}(m)
= 2(m{+}1)/(2m{+}1) = I_m/I_{m+1}$
propagates inductively from $S^{(0)} = 1/I_0$,
yielding $\operatorname{val}(m) = 1/I_m =
2\,\Gamma(m{+}1)/(\sqrt{\pi}\,\Gamma(m{+}\half))$
for all non-negative integers~$m$.
Under the same assumption, the identity extends to all real
$m > -\frac{1}{2}$ by analytic continuation.
\end{proof}

\begin{remark}
For $m = 0$ and $m = 1$, Steps~1--3 give a complete algebraic proof
independent of Conjecture~\ref{conj:ratio}:
$S^{(0)} = 2/\pi$ and $S^{(1)} = 4/\pi$.
For $m \ge 2$, the proof is conditional on the ratio conjecture.
No polynomial-coefficient shift operator $L_m$ mapping $\mathcal{S}_m$
to $\mathcal{S}_{m+1}$ has been found for general~$m$.
\end{remark}

\begin{corollary}[{\normalfont(conditional on Conjecture~\ref{conj:ratio})}]\label{cor:special}
Assuming Conjecture~\ref{conj:ratio}, the following specializations hold.
\begin{enumerate}
\item[\textup{(i)}]
  Integer $m \ge 0$:
  $\operatorname{val}(m) = 2^{2m+1}/\bigl(\pi\binom{2m}{m}\bigr)$.
\item[\textup{(ii)}]
  Half-integer $m = k{+}\half$, $k \ge 0$:
  $\operatorname{val}(m)$ is an exact rational (Wallis product partial convergent).
\item[\textup{(iii)}]
  $\operatorname{val}(m) \sim \sqrt{4m/\pi}$
  as $m \to \infty$ by Stirling's formula.
\end{enumerate}
\end{corollary}

\begin{proof}
Part~\textup{(i)} is the integer case of Theorem~\ref{thm:gamma}.
For part~\textup{(ii)}, substitute $m=k{+}\half$ into the Gamma quotient and
use $\Gamma(k{+}1)=k!$. Part~\textup{(iii)} is the corresponding Stirling
asymptotic for the same quotient.
\end{proof}

%======================================================================
\section{Casoratian Series Identity for \texorpdfstring{$\pi/4$}{pi/4}}%
\label{sec:casoratian}
%======================================================================

We give a Casoratian derivation of the series identity for~$\pi/4$.
The telescoping argument is internal to the $m=1$ recurrence, while
the identification of the continued-fraction limit as~$4/\pi$ comes
independently from the Euler/Wallis analysis of the previous section.

\subsection{Convergence}

\begin{proposition}\label{prop:conv}
The CF~\eqref{eq:pcf} with
$\alpha(n) = -n(2n{-}3)$, $\beta(n) = 3n{+}1$ converges.
\end{proposition}

\begin{proof}
By the \'Sleszy\'nski--Pringsheim theorem
\citep[Theorem~4.35]{lorentzen2008},
the CF converges if
$|\alpha(n)| \leq \beta(n)\,\beta(n{-}1)$ for $n \geq 2$.
We have
$\beta(n)\,\beta(n{-}1) - |\alpha(n)|
= (3n{+}1)(3n{-}2) - n(2n{-}3) = 7n^2 - 2$,
which is strictly positive for $n \geq 1$.
\end{proof}

\subsection{Casoratian recursion}

Let $h_n = n^2{+}3n{+}1$ (the polynomial solution from
Theorem~\ref{thm:conv1}) and $g_n = q_n/\dff{2n{-}1}$,
both satisfying the reduced recurrence~\eqref{eq:ured}.

\begin{lemma}[Casoratian recursion]\label{lem:cas_rec}
The Casoratian $W_n = h_n\,g_{n-1} - h_{n-1}\,g_n$ satisfies
$W_n = \frac{n}{2n{-}1}\,W_{n-1}$ for $n \geq 2$.
\end{lemma}

\begin{proof}
Substituting the recurrence~\eqref{eq:ured} for $h_n$ and $g_n$:
\begin{align*}
  W_n
  &= \frac{(3n{+}1)h_{n-1} - n\,h_{n-2}}{2n{-}1}\,g_{n-1}
    - h_{n-1}\,\frac{(3n{+}1)g_{n-1} - n\,g_{n-2}}{2n{-}1} \\
  &= \frac{n}{2n{-}1}(h_{n-1}\,g_{n-2} - h_{n-2}\,g_{n-1})
   = \frac{n}{2n{-}1}\,W_{n-1}.\qedhere
\end{align*}
\end{proof}

\begin{theorem}[Casoratian closed form]\label{thm:cas}
For all $n \geq 1$,
$W_n = n!/\dff{2n{-}1}$.
\end{theorem}

\begin{proof}
Base: $W_1 = 5\cdot 1 - 1\cdot 4 = 1 = 1!/1!!$.
Induction:
$W_n = \frac{n}{2n{-}1}\cdot\frac{(n{-}1)!}{\dff{2n{-}3}}
= \frac{n!}{\dff{2n{-}1}}$.
\end{proof}

\subsection{Main series identity}

\begin{theorem}\label{thm:4overpi}
Let $S$ be the CF with $\alpha(n) = -n(2n{-}3)$,
$\beta(n) = 3n{+}1$.
\begin{enumerate}
\item[\textup{(i)}]
  $S$ converges, with $|S_n - S| = O(2^{-n}/n^{7/2})$.
\item[\textup{(ii)}]
  $p_n = \dff{2n{-}1}\,(n^2{+}3n{+}1)$.
\item[\textup{(iii)}]
  $S = 4/\pi$, equivalently
\begin{equation}\label{eq:series}
  \frac{\pi}{4}
  = 1 - \sum_{k=1}^{\infty}
    \frac{k!}{\dff{2k{-}1}\,(k^2{+}3k{+}1)(k^2{+}k{-}1)}
  = 1 - \sum_{k=1}^{\infty}
    \frac{2^k}{\binom{2k}{k}(k^2{+}3k{+}1)(k^2{+}k{-}1)}.
\end{equation}
\end{enumerate}
\end{theorem}

\begin{proof}[Proof of \textup{(iii)}]
Part~\textup{(ii)} is exactly the closed numerator formula~\eqref{eq:pn_m1}
from Theorem~\ref{thm:conv1}. Part~\textup{(i)} combines
Proposition~\ref{prop:conv} with the tail estimate proved immediately below.
It remains to establish \textup{(iii)}. The ratio $g_n/h_n$ telescopes:
\[
  \frac{g_n}{h_n} - \frac{g_{n-1}}{h_{n-1}}
  = \frac{-W_n}{h_n\,h_{n-1}}.
\]
Summing from $k = 1$ to~$n$ and using $g_0/h_0 = 1$:
\[
  \frac{g_n}{h_n}
  = 1 - \sum_{k=1}^{n}
    \frac{k!}{\dff{2k{-}1}\,(k^2{+}3k{+}1)(k^2{+}k{-}1)}.
\]
As $n\to\infty$, the convergents $p_n/q_n = h_n/g_n$ tend to the
continued-fraction value $S$ by Proposition~\ref{prop:conv}.
The value $S = 4/\pi$ is established independently by
Lemma~\ref{lem:shift} and Step~3 of the proof of
Theorem~\ref{thm:gamma}: the raising operator $L$ maps
$P^{(0)},Q^{(0)}$ to scalar multiples of $P^{(1)},Q^{(1)}$, so
$S^{(1)} = 2\,S^{(0)} = 4/\pi$ by Theorem~\ref{thm:pi0}, with no
appeal to the present section. Therefore $g_n/h_n \to \pi/4$, and
passing to the limit in the telescoping identity yields
\eqref{eq:series}. The alternative form uses
$k!/\dff{2k{-}1} = 2^k/\binom{2k}{k}$
\citep[see][\S1.2.6]{knuth1997}.
\end{proof}

\begin{proof}[Proof of convergence rate \textup{(i)}]
By Stirling, $\binom{2k}{k} \sim 4^k/\sqrt{\pi k}$, so the
general term satisfies
$t_k \sim \sqrt{\pi k}/(2^k k^4) = O(2^{-k}/k^{7/2})$.
Since the tail is dominated by its first term,
$|S_N - \pi/4| = O(2^{-N}/N^{7/2})$.
At $N = 35$ the error drops below~$10^{-16}$.
\end{proof}

%======================================================================
\section{Gauss CF Non-Membership}\label{sec:gauss}
%======================================================================

\subsection{Factorisation over \texorpdfstring{$\mathbb{Q}(\sqrt{5})$}{Q(sqrt5)}}

Set $\varphi = (1{+}\sqrt{5})/2$, $\psi = (1{-}\sqrt{5})/2$.
Since $(k{+}\varphi)(k{+}\psi) = k^2{+}k{-}1$
and $(k{+}2{+}\varphi)(k{+}2{+}\psi) = k^2{+}5k{+}5$,
the ratio of consecutive series terms in~\eqref{eq:series} is
\begin{equation}\label{eq:ratio_hyp}
  \frac{t_{k+1}}{t_k}
  = \frac{(k{+}1)(k{+}\varphi)(k{+}\psi)}
         {2(k{+}\half)(k{+}2{+}\varphi)(k{+}2{+}\psi)}.
\end{equation}

\subsection{Pochhammer obstruction}

A standard $\hyp{p}{q}{\cdots}{\cdots}{z}$ has term ratio equal to
a product of linear factors $(k+a_i)/(k+b_j)$ with
$a_i,b_j\in\mathbb{Q}$.
Here the irreducible quadratics $k^2{+}k{-}1$ and $k^2{+}5k{+}5$
(both of discriminant~$5$) force the parameters into
$\mathbb{Q}(\sqrt5)\setminus\mathbb{Q}$.

\subsection{Gauss CF non-membership}

\begin{proposition}\label{prop:gauss}
The CF does not arise from any ${}_2F_1$ ratio via the Gauss
continued-fraction theorem.
\end{proposition}

\begin{proof}
The Gauss CF for
$\hyp21{A{+}1,\,B}{C{+}1}{z}\big/\hyp21{A,B}{C}{z}$
has unit-denominator coefficients $t_k$ determined by
$(A,B,C)\in\mathbb{Q}^3$.
Solving $t_1=c_1$, $t_2=c_2$, $t_3=c_3$ (where the $c_k$ are the
equivalence-transformed coefficients of our CF) yields a degree-6
polynomial system with six solutions; four are complex and two are
real with $C = -3.2590\ldots$ (certified in Arb at $50$-digit precision).
For these,
$|t_4 - c_4| = 5.900\ldots$ (certified), so no match exists.
The polynomial system was solved using \texttt{SymPy} and all
coefficient values verified at $50$-digit precision in Arb.
\end{proof}

\subsection{\texorpdfstring{${}_2F_3$}{2F3} identification}

\begin{remark}[${}_2F_3$ identification]\label{rem:2f3}
The term ratio~\eqref{eq:ratio_hyp} formally identifies the series as
\[
  {}_2F_3\!\left(\varphi,\psi;\;\half,\;2{+}\varphi,\;2{+}\psi;\;
  \half\right)
\]
over $\mathbb{Q}(\sqrt5)$: the upper parameters
$\varphi = (1{+}\sqrt5)/2$ and $\psi = (1{-}\sqrt5)/2$ match the
factored Pochhammer shifts in~\eqref{eq:ratio_hyp}, while the lower
parameters $\frac{1}{2}$, $2{+}\varphi$, $2{+}\psi$ correspond to the
denominator factors $(k{+}\frac{1}{2})$, $(k{+}2{+}\varphi)$,
$(k{+}2{+}\psi)$.
Since both the parameters and the function class depend on the
extension $\mathbb{Q}(\sqrt5)/\mathbb{Q}$, the series does not
reduce to any standard ${}_pF_q$ with rational parameters.
We state this identification as a structural observation rather than
a theorem; a full verification that the Pochhammer products agree
term-by-term for all~$k$ is a routine but unilluminating exercise in
shifted-factorial bookkeeping.
\end{remark}

%======================================================================
\section{Euler CF Duality}\label{sec:euler}
%======================================================================

The $m=0$ CF decomposes as
$S^{(0)} = \beta(0) + \alpha(1)/T = 1 - 1/T$
where~$T$ is the inner tail starting at $n = 1$.
Since $S^{(0)} = 2/\pi$ (Theorem~\ref{thm:pi0}):

\begin{proposition}\label{prop:euler}
$T = \pi/(\pi - 2)$, equivalently
$T/(T{-}1) = \pi/2$.
\end{proposition}

\begin{proof}
$T = 1/(1 - S^{(0)}) = 1/(1 - 2/\pi) = \pi/(\pi - 2)$.
The identity $T/(T{-}1) = \pi/2$ then follows from
$T - 1 = (2)/(\pi - 2)$.
\end{proof}

\begin{remark}
This connects our PCF to the classical identity~\eqref{eq:pi-id}, which comes
from $\arcsin(1)$ \citep[\S4.24.2]{dlmf}.
The partial numerators of the $m=0$ and $m=1$ families differ
by a linear term:
$\alpha_0(n) - \alpha_1(n) = -n(2n{-}1) + n(2n{-}3) = -2n$.
\end{remark}

%======================================================================
\section{Comparison with Brouncker's CF}\label{sec:brouncker}
%======================================================================

Lord Brouncker's classical CF for $4/\pi$ is
\[
  \frac{4}{\pi}
  = 1 + \cfrac{1^2}{2 + \cfrac{3^2}{2 + \cfrac{5^2}{2 + \cdots}}},
\]
with $\alpha^B(n) = (2n{-}1)^2$ and constant $\beta^B(n) = 2$.
Both CFs converge to~$4/\pi$ but differ structurally:

\begin{center}
\begin{tabular}{@{}lcc@{}}
  \toprule
  & \textbf{Brouncker} & \textbf{Novel ($m{=}1$)} \\
  \midrule
  $\alpha(n)$ & $(2n{-}1)^2$ & $-n(2n{-}3)$ \\
  $\beta(n)$ & $2$ (constant) & $3n{+}1$ (linear) \\
  $p_n$ & $(2n{+}1)!!$ & $\dff{2n{-}1}(n^2{+}3n{+}1)$ \\
  Series & Leibniz $\sum (-1)^k/(2k{+}1)$ &
    Eq.~\eqref{eq:series} \\
  Rate & $O(1/n)$ & $O(2^{-n})$ \\
  \bottomrule
\end{tabular}
\end{center}

\begin{center}
\small
\begin{tabular}{@{}rrrcrrc@{}}
  \toprule
  & \multicolumn{3}{c}{\textbf{Brouncker}} &
    \multicolumn{3}{c}{\textbf{Novel}} \\
  \cmidrule(lr){2-4}\cmidrule(lr){5-7}
  $n$ & $p_n$ & $q_n$ & $S_n$ & $p_n$ & $q_n$ & $S_n$ \\
  \midrule
  0 & 1 & 1 & 1.000000 & 1 & 1 & 1.000000 \\
  1 & 3 & 2 & 1.500000 & 5 & 4 & 1.250000 \\
  2 & 15 & 13 & 1.153846 & 33 & 26 & 1.269231 \\
  3 & 105 & 76 & 1.381579 & 285 & 224 & 1.272321 \\
  4 & 945 & 789 & 1.197719 & 3045 & 2392 & 1.272993 \\
  5 & 10395 & 7734 & 1.344065 & 38745 & 30432 & 1.273166 \\
  8 & 34459425 & 28018665 & 1.229874 &
    180405225 & 141690240 & 1.273237 \\
  \bottomrule
  \multicolumn{7}{r}{$4/\pi \approx 1.273\,239\,545$}
\end{tabular}
\end{center}

\noindent
At $n = 8$, Brouncker's error is $\sim\!4\times 10^{-2}$ while the
novel CF error is $\sim\!3\times 10^{-6}$, a factor of~$10^4$.
The two CFs are not related by any equivalence transformation,
contraction, or M\"obius substitution.

%======================================================================
\section{Convergence Theory}\label{sec:conv}
%======================================================================

\subsection{Worpitzky condition}

\begin{lemma}\label{lem:worpitzky}
For both families, the normalised partial numerators
$|\tilde\alpha_n| = |\alpha_n|/(\beta_n\,\beta_{n-1})$ satisfy
$|\tilde\alpha_n| < 1/4$ for all sufficiently large~$n$.
\end{lemma}

\begin{proof}
\textbf{Log family.}
$|\tilde\alpha_n| \to k/(k{+}1)^2 < 1/4$ for $k > 1$
(since $(k{-}1)^2 > 0$).

\textbf{Pi family.}
$|\tilde\alpha_n| = n(2n{-}1)/((3n{+}1)(3n{-}2)) \to 2/9 < 1/4$.
At $n = 1$: $|\tilde\alpha_1| = 1/4$ (borderline);
strict inequality holds for $n \ge 2$, and convergence at
equality at finitely many~$n$ follows from the classical
extension \citep[\S4.2]{lorentzen2008}.
\end{proof}

\subsection{Bracketing}

\begin{lemma}\label{lem:bracket}
If $\alpha_n < 0$ and $\beta_n > 0$ for all $n \ge 1$, and the CF
converges to~$L$, then
$\min(C_N,C_{N+1}) \le L \le \max(C_N,C_{N+1})$.
\end{lemma}

\begin{proof}
The determinant $p_n q_{n-1} - p_{n-1} q_n
= \prod_{j=1}^n (-\alpha_j) > 0$
shows $C_n - C_{n-1}$ alternates in sign.
See \citet[Theorem~4.35]{lorentzen2008}.
\end{proof}

%======================================================================
\section{Meta-Family}\label{sec:meta}
%======================================================================

\begin{remark}[Unified template]
The template
$\alpha(n) = -(\alpha_0 n^2 + \beta_0 n)$,
$\beta(n) = \gamma_0 n + \delta_0$
contains both the Logarithmic Ladder
($\beta_0 = 0$, $\delta_0 = \alpha_0$, $\gamma_0 = \alpha_0 + 1$)
and the Pi Family
($\alpha_0 = 2$, $\gamma_0 = 3$, $\delta_0 = 1$)
as special cases.
Two additional sub-families,
$(\alpha_0, \beta_0, \gamma_0, \delta_0) = (1,0,4,2) \to 2/\ln 3$
and $(4,0,8,4) \to 4/\ln 3$,
were found and certified at $10\,000$-bit precision with depth $N = 5000$;
both are instances of the Logarithmic Ladder at~$k = 3$
(and its equivalence-transformed double).
\end{remark}

%======================================================================
\section{Quadratic GCF Constants}\label{sec:vquad}
%======================================================================

Define $V(A,B,C) = 1 + \mathbf{K}_{n \ge 1}\; 1/(An^2 + Bn + C)$.

\begin{lemma}[Denominator growth]\label{lem:Qgrowth}
Let $A \ge 1$, $C \ge 1$, and $An^2{+}Bn{+}C > 0$ for $n \ge 1$.
Then the convergent denominators
$Q_n = (An^2{+}Bn{+}C)\,Q_{n-1} + Q_{n-2}$
satisfy $\log Q_n \ge n\log n - O(n)$.
\end{lemma}

\begin{proof}
Since $Q_{n-2} \ge 0$ and $Q_0 = 1$, $Q_1 = A{+}B{+}C \ge 1$,
we have $Q_n \ge (An^2{+}Bn{+}C)\,Q_{n-1}$ for all $n \ge 1$.
Iterating:
\[
  Q_n \ge \prod_{j=1}^{n}(Aj^2{+}Bj{+}C).
\]
For $j \ge j_0$ with $j_0$ depending only on $A,B,C$, we have
$Aj^2{+}Bj{+}C \ge j^2/2$, so
\[
  \log Q_n
  \ge \sum_{j=j_0}^{n} \log(j^2/2) + O(1)
  = 2\sum_{j=j_0}^{n}\log j - (n{-}j_0{+}1)\log 2 + O(1).
\]
By Stirling, $\sum_{j=1}^{n}\log j = n\log n - n + O(\log n)$,
giving $\log Q_n \ge 2n\log n - O(n)$, and in particular
$\log Q_n \ge n\log n$ for all sufficiently large~$n$.
\end{proof}

\begin{proposition}\label{prop:irrat}
For all $(A,B,C)$ with $A \ge 1$, $C \ge 1$, and
$An^2 + Bn + C > 0$ for $n \ge 1$, $V(A,B,C)$ is irrational.
\end{proposition}

\begin{proof}
The Wronskian identity
$|P_n Q_{n-1} - P_{n-1} Q_n| = 1$ holds for all~$n$
(since all partial numerators equal~$1$).
By Lemma~\ref{lem:Qgrowth}, $Q_n$ grows superexponentially,
so $|V - P_n/Q_n| = 1/(Q_n Q_{n+1}) = o(Q_n^{-1})$.
Euler's irrationality criterion then gives irrationality of~$V$.
\end{proof}

Our systematic scan found \textbf{482 distinct} such constants.
All values in this paragraph were first computed at $300$ decimal digits with
forward depth $N=2000$, and the final Arb enclosures are recorded in
Appendix~\ref{sec:cert}.
A PSLQ \citep{ferguson1999} integer relation search at 300-digit precision over 15
standard constants
$\{\pi, \pi^2, e, \ln 2, \ln 3, \gamma, G, \zeta(3), \zeta(5),
\sqrt{2}, \sqrt{3}, \sqrt{5}, \varphi, \sqrt{\pi}, \pi\ln 2\}$
finds no relation for any $V$-constant, and none is algebraic of
degree~$\le 8$ with coefficients~$\le 10^5$.
Pairwise cross-relation tests on 91 representative constants
(435~pairs, PSLQ with coefficients up to~500) and triple-relation
tests (455~triples, coefficients up to~200) all return negative.
A targeted sweep against 588 expressions built from Gamma ratios
and Pi Family values also produces no matches.

These computations have a mathematical role beyond data collection.
Proposition~\ref{prop:irrat} gives a uniform irrationality theorem for the
quadratic family, and the negative PSLQ evidence indicates that the 482 values
are not simple reparameterisations of the standard constants already in the
search basis.

The only structural relation found is the \emph{shift symmetry}
$V(1,B,C) - V(1,-B,C) \in \mathbb{Z}$ for all tested $B, C$,
following from the palindromic structure of the GCF denominators.
This gives strong (but non-rigorous) evidence of mutual algebraic
independence.

\begin{remark}
An earlier PSLQ search returned $[0, 1, -2, 1]$ relative to
$\{V, \sqrt{5}, \varphi, 1\}$; this is the trivial identity
$\sqrt{5} - 2\varphi + 1 = 0$ with zero coefficient on~$V$.
\end{remark}

%======================================================================
\appendix
\section{Numerical Certification}\label{sec:cert}
%======================================================================

All results are certified using Arb ball arithmetic
(\texttt{python-flint}~0.8.0, $10\,000$-bit precision)
at depth $N = 5000$.
By Lemma~\ref{lem:bracket}, the limit lies in
$[\min(C_N, C_{N-1}),\, \max(C_N, C_{N-1})]$.
Table~\ref{tab:cert} summarises the main certified enclosures.
These computations serve as certification of the proved formulas and of the
catalogue values recorded in Section~\ref{sec:vquad}.

\begin{table}[ht]
\centering
\caption{Arb-certified intervals (depth~5000)}\label{tab:cert}
\begin{tabular}{@{}lccc@{}}
\toprule
PCF & Target & Bracket width & Digits \\
\midrule
$\alpha{=}{-}2n^2,\;\beta{=}3n{+}2$
  & $1/\ln 2$    & $<10^{-1509}$ & 1509 \\
$\alpha{=}{-}3n^2,\;\beta{=}4n{+}3$
  & $1/\ln(3/2)$ & $<10^{-2265}$ & 2265 \\
$\alpha{=}{-}n(2n{-}1),\;\beta{=}3n{+}1$
  & $2/\pi$      & $<10^{-1508}$ & 1508 \\
$\alpha{=}{-}n(2n{-}3),\;\beta{=}3n{+}1$
  & $4/\pi$      & $<10^{-1518}$ & 1518 \\
\bottomrule
\end{tabular}
\end{table}

For the $m = 1$ CF, backward recurrence at depth~500 with
80-digit arithmetic confirms
\[
  S = 1.27323954473516268615107010698011489627567716592365158998\ldots
\]
matching $4/\pi$ to all 60 digits.

%======================================================================
\section{Discovery Methodology}\label{sec:method}
%======================================================================

\begin{enumerate}
\item \textbf{MITM sweep.}
  Fix $\beta(n) = dn{+}e$; enumerate
  $\alpha(n) = pn^2{+}qn{+}r$ with $|p|,|q|,|r| \le 5$;
  match PCF values against a database of 200 known constants.
\item \textbf{Modular signature sweep.}
  Vary $d \in \{2,\ldots,8\}$, $e \in \{1,\ldots,8\}$.
  The $d{=}3$ family produced both $\pi$- and $\ln$-hits.
\item \textbf{Parametric generalisation.}
  Sweep $c$ in $\alpha(n) = -n(2n{-}c)$ and $k$ in
  $\alpha(n) = -kn^2$, discovering the full infinite families.
\end{enumerate}

%======================================================================
\section{Experimental Cubic and Quartic Clausen Witnesses}\label{sec:exp-cq}
%======================================================================

The proved families above motivated a parallel high-precision search in nearby
cubic J-fractions. The numerically strongest cubic case is the one-parameter
family
\[
  B(\beta) = \beta + \mathbf{K}_{n\ge1}\frac{2n^3}{2n+\beta},
  \qquad \beta \in \mathbb{Q}_{>0}.
\]
At the anchor point $\beta=\tfrac{2}{3}$,
\[
\begin{aligned}
  B_{\mathrm{cubic}}
    &= \operatorname{eval\_cf}(2,3,2,2/3)\\
    &= \frac23 + \mathbf{K}_{n\ge1}\frac{2n^3}{2n+\frac23}
     = 1.110030230754775938693\ldots.
\end{aligned}
\]
With $\theta = \arcsin(3\sqrt{3}/8)$, this matches
\[
  {}_2F_1\!\left(u,1-u;\tfrac{3}{2};\tfrac{27}{64}\right)^2
  = \left(\frac{\sin(\lambda\theta)}{\lambda\sin\theta}\right)^2,
  \qquad \lambda = 1-2u,
\]
for $u = 0.18852361548552202250\ldots$, with hypergeometric residual below
$10^{-341}$ in independent \texttt{mpmath} checks at 350-digit precision.

A companion evaluation at $\beta=\tfrac{1}{3}$ gives
\[
  B\!\left(\tfrac{1}{3}\right)
  = 0.805664740414447796515\ldots
  = {}_2F_1\!\left(u',1-u';\tfrac{3}{2};\tfrac{27}{64}\right)^2,
\]
with $u' = -0.24589471253854869965\ldots$ on the $\lambda>1$ sheet of the same
cubic modular point $z=27/64$.

\begin{remark}
Exhaustive integer-relation searches at 300-digit precision, including \texttt{findpoly} to degree~20, PSLQ over mixed bases involving $\sqrt{3}$ and the elliptic periods $K(k),K(k')$ with $k=3\sqrt{3}/8$, and pairwise combination tests over all five computed $\lambda$-values, find no algebraic structure for the deformation parameter $\lambda(\beta)$.
\end{remark}

\begin{remark}[Structural map of the point $27/64$]
The point $z=27/64$ is the ramification point of the cubic Chebyshev resolvent
$4x^3 - 3x + \frac{8}{3\sqrt3}\sqrt{z}=0$,
whose discriminant $\Delta_x = -16(64z-27)$ vanishes exactly at $z = 27/64$.
In the variable $t = \sqrt{1 - 64z/27}$, the pulled-back Gauss equation for
${}_2F_1(u,1{-}u;\frac{3}{2};z)$ acquires a regular singularity at $t=0$.
The integers $27 = 3^3$ and $64 = 4^3$ arise from the cubic Pochhammer
combinatorics $(\frac{1}{3})_n(\frac{2}{3})_n = (3n)!/(27^n n!)$ and the
Chebyshev normalization, respectively.

The real branch persists up to a numerically observed endpoint
$\beta_{\mathrm{crit}} \approx 0.7478702596$ where
$\lambda(\beta_{\mathrm{crit}}) = 0$ and the value saturates at
$(\theta/\sin\theta)^2 \approx 1.18466547096$.
The associated $j$-invariant $j(27/64) = 29704593673/15968016$ is rational
but not an algebraic integer, so $z = 27/64$ is not a classical CM singular
modulus; the point behaves as a \emph{modular shadow} rather than a genuine
cusp of the continued fraction.
\end{remark}

A preliminary scan at the quartic point $z=\tfrac{1}{4}$ for the family
$\operatorname{eval\_cf}(1,4,7,\beta)$ shows a weak shadow at five-digit
precision; this is recorded as a direction for future investigation in
\S\ref{sec:open}.

%======================================================================
\section{Open Questions}\label{sec:open-cubic}\label{sec:open}
%======================================================================

The following questions arise directly from the proved results and
numerical evidence.

\begin{enumerate}
\item \textbf{Prove Conjecture~\ref{conj:ratio}.}
The ratio formula $\operatorname{val}(m{+}1) = \operatorname{val}(m)\cdot 2(m{+}1)/(2m{+}1)$
is verified numerically to at least 500 digits for $m = 0,1,\ldots,10$
and half-integer values (Section~\ref{sec:pi}).
A proof would upgrade all conditional results in this paper
(Conditional Theorem~\ref{thm:gamma}, Corollaries~\ref{cor:pi}
and~\ref{cor:special}) to unconditional theorems and complete the
continuous Gamma-function identity for all~$m$.
This is the primary open problem of the paper.

\item \textbf{Proof of the Clausen-square identity.}
Prove that $B(2/3) = [{}_{2}F_1(u,1-u;\tfrac{3}{2};\tfrac{27}{64})]^2$
for $u = 0.18852361548552202250\ldots$ as an exact identity,
not merely a high-precision numerical coincidence.
The classical Clausen identity provides a structural scaffold;
what is missing is a direct derivation connecting
$\operatorname{eval\_cf}(2,3,2,\beta)$ to the hypergeometric value at this argument.

\item \textbf{Nature of the deformation parameter $\lambda(\beta)$.}
For each $\beta$ in the real branch $[0, \beta_{\mathrm{crit}}]$,
the map $\beta \mapsto \lambda(\beta)$ is smooth and monotone decreasing.
Integer-relation searches at 300-digit precision find no algebraic
relation for $\lambda(\beta)$ at degree $\le 20$ with coefficient
height $\le 10^5$, no relation involving $\sqrt{3}$ or the elliptic
periods $K(k), K(k')$ for $k = 3\sqrt{3}/8$, and no relation
of the form $\lambda(\beta)\theta/\pi \in \mathbb{Q}$.
Whether $\lambda(\beta)$ is transcendental for all rational $\beta$,
or algebraic at a sparse set of special values, remains open.

\item \textbf{Non-classical modular status and the quadratic catalogue.}
The elliptic modulus $k = 3\sqrt{3}/8$ satisfies
$k^3 + (k')^3 = 0.71359\ldots \ne 1$,
placing $z = 27/64$ strictly outside the Borwein signature-3
singular locus.
The associated nome $q = 0.034194\ldots$ matches no entry
$q_r = e^{-\pi\sqrt{r}}$ for $r \le 36$.
What modular or arithmetic object, if any, organizes the
selection of $z = 27/64$ by the continued-fraction family
$\operatorname{eval\_cf}(2,3,2,\beta)$?
In the quadratic catalogue, the constant $V(3,1,1)$ is characterized
in a companion paper \citep{papanokechi2026spectral} as a
confluent-Heun / Stokes obstruction within the spectral classification
of polynomial continued fractions; the precise relationship between
these two obstruction mechanisms remains to be clarified.

\item \textbf{The real-branch endpoint.}
The map $\beta \mapsto V(\beta)$ appears to have a critical point
near $\beta_{\mathrm{crit}} \approx 0.74787025960$, where
$V(\beta_{\mathrm{crit}}) \approx (\theta/\sin\theta)^2$,
the $\lambda \to 0$ limit of the trigonometric form.
Does this endpoint admit a closed form,
and does the family admit a natural complex or hyperbolic
continuation beyond it?

\item \textbf{Higher-degree analogues.}
Does an analogous Clausen-square collapse occur for
$\operatorname{eval\_cf}(p, d, \alpha, \beta)$ with $d \ge 4$?
A preliminary scan at the quartic point $z = 1/4$
for $\operatorname{eval\_cf}(1,4,7,\beta)$ shows a weak match
at five-digit precision; this deserves systematic
high-precision investigation.

\item \textbf{Casoratian telescoping for higher-rank recurrences.}
Determine whether the Casoratian telescoping method of
Section~\ref{sec:casoratian} extends to higher-rank recurrences
($d \geq 2$) to yield closed-form evaluations beyond the logarithmic
and $\pi$ sectors.

\end{enumerate}

%======================================================================
\section{Exact Convergents}\label{app:conv}
%======================================================================

\textbf{Log ($k{=}2$):} $p_n = (n{+}1)!\cdot 2^{n+1}$.

\smallskip
\begin{center}
\begin{tabular}{@{}rrrl@{}}
\toprule
$n$ & $p_n$ & $q_n$ & $C_n$ \\
\midrule
0 & 2 & 1 & 2.000000 \\
1 & 8 & 5 & 1.600000 \\
2 & 48 & 32 & 1.500000 \\
3 & 384 & 262 & 1.465649 \\
4 & 3840 & 2644 & 1.452345 \\
5 & 46080 & 31832 & 1.447584 \\
\bottomrule
\end{tabular}
\end{center}
\smallskip\noindent Limit: $1/\ln 2 = 1.442695\ldots$

\medskip
\textbf{Pi ($m{=}0$):} $p_n = (2n{+}1)!!$.

\smallskip
\begin{center}
\begin{tabular}{@{}rrrl@{}}
\toprule
$n$ & $p_n$ & $q_n$ & $C_n$ \\
\midrule
0 & 1 & 1 & 1.000000 \\
1 & 3 & 4 & 0.750000 \\
2 & 15 & 22 & 0.681818 \\
3 & 105 & 160 & 0.656250 \\
4 & 945 & 1464 & 0.645492 \\
5 & 10395 & 16224 & 0.640688 \\
\bottomrule
\end{tabular}
\end{center}
\smallskip\noindent Limit: $2/\pi = 0.636620\ldots$

\medskip
\textbf{Pi ($m{=}1$):}
$p_n = \dff{2n{-}1}\,(n^2{+}3n{+}1)$.

\smallskip
\begin{center}
\begin{tabular}{@{}rrl@{}}
\toprule
$n$ & $p_n$ & $q_n$ \\
\midrule
0 & 1 & 1 \\
1 & 5 & 4 \\
2 & 33 & 26 \\
3 & 285 & 224 \\
4 & 3\,045 & 2\,392 \\
5 & 38\,745 & 30\,432 \\
6 & 571\,725 & 449\,040 \\
7 & 9\,594\,585 & 7\,535\,616 \\
8 & 180\,405\,225 & 141\,690\,240 \\
\bottomrule
\end{tabular}
\end{center}
\smallskip\noindent Limit: $4/\pi = 1.273239\ldots$

%======================================================================
\section{Reproducibility}\label{app:repro}
%======================================================================

Python~3.12, \texttt{mpmath}~1.4.1,
\texttt{python-flint}~0.8.0 (Arb),
\texttt{sympy}~1.14.0.
Forward recurrence~\eqref{eq:rec} with
$p_{-1}{=}1$, $p_0{=}\beta(0)$, $q_{-1}{=}0$, $q_0{=}1$.
Runtime: ${<}\,2$~hours on a standard desktop.
The symbolic verification of Conjecture~\ref{conj:Rnm}
(for $m=0,\ldots,12$ and $n\le 20$) is performed by
\texttt{conj43\_check.py} in the repository.

%======================================================================
\section*{Acknowledgements}
%======================================================================

Inspired by the Ramanujan Machine \citep{raayoni2021,elimelech2024}.
Computations performed with \texttt{mpmath} and \texttt{python-flint} (Arb).

% EDITORIAL NOTE: AI DISCLOSURE UPDATED
\paragraph{Computational Tools.}
This research was conducted using the SIARC
(Self-Iterating Analytic Relay Chain) v6.5 methodology, employing
AI language models as research assistants throughout the discovery
and verification workflow.
Specific tools used include Anthropic Claude (claude.ai) and GitHub
Copilot (GPT-4 backend) for hypothesis generation, symbolic computation
review, proof scaffolding, and manuscript preparation.
All mathematical claims, numerical values, and proof steps were
independently verified by the author using certified Arb ball arithmetic
(\texttt{python-flint}~0.8.0), \texttt{mpmath}~1.4.1, and
\texttt{SymPy}~1.14.0, with no mathematical content sourced from AI
output without independent verification.

The epistemic governance framework governing AI use in this research is
documented in:
(i)~the AEAL (Agent Epistemic Accountability Layer)
framework~\citep{papanokechi2026aeal};
(ii)~the ZTEK (Zero Trust Epistemic Kernel)
protocol~\citep{papanokechi2026ztek}; and
(iii)~the AI Behavior Science founding
paper~\citep{papanokechi2026aibs}.
This disclosure follows \textit{Experimental Mathematics} editorial
policy on AI-assisted research.

\subsection*{Data Availability}

All computation scripts and Arb-certified output files are available at
\url{https://github.com/papanokechi/pcf-research}.

%======================================================================
\begin{thebibliography}{12}
%======================================================================

\bibitem[Ap\'ery(1979)]{apery1979}
R.~Ap\'ery,
\textit{Irrationalit\'e de $\zeta(2)$ et $\zeta(3)$},
Ast\'erisque \textbf{61} (1979), 11--13.

\bibitem[Cuyt et~al.(2008)]{cuyt2008}
A.~Cuyt, V.\,B.~Petersen, B.~Verdonk, H.~Waadeland,
and W.\,B.~Jones,
\emph{Handbook of Continued Fractions for Special Functions},
Springer, 2008.

\bibitem[DLMF(2024)]{dlmf}
F.\,W.\,J.~Olver, A.\,B.~Olde~Daalhuis, D.\,W.~Lozier,
et~al.\ (eds.),
\emph{NIST Digital Library of Mathematical Functions},
Release~1.2.3, \url{https://dlmf.nist.gov/}, 2024.

\bibitem[Elimelech et~al.(2024)]{elimelech2024}
D.~Elimelech et~al.,
\textit{Algorithm-assisted discovery of an intrinsic order among
mathematical constants},
Proc.\ Natl.\ Acad.\ Sci.\ USA \textbf{121}(25) (2024),
e2321440121.

\bibitem[Ferguson et~al.(1999)]{ferguson1999}
H.\,R.\,P.~Ferguson, D.\,H.~Bailey, and S.~Arwade,
\textit{Analysis of PSLQ, an integer relation finding algorithm},
Math.\ Comp.\ \textbf{68} (1999), 351--369.

\bibitem[Knuth(1997)]{knuth1997}
D.\,E.~Knuth,
\emph{The Art of Computer Programming},
vol.~1, 3rd ed., Addison-Wesley, 1997.

\bibitem[Lorentzen and Waadeland(2008)]{lorentzen2008}
L.~Lorentzen and H.~Waadeland,
\emph{Continued Fractions with Applications},
2nd ed., Atlantis Press, 2008.

\bibitem[{OEIS Foundation}(2024)]{oeis_A028387}
{OEIS Foundation},
\textit{A028387: $n + (n{+}1)^2$},
The On-Line Encyclopedia of Integer Sequences,
\url{https://oeis.org/A028387}.

\bibitem[Papanokechi(2026)]{papanokechi2026spectral}
Papanokechi,
\textit{Spectral Classes of Polynomial Continued Fractions:
A Hypergeometric Framework and Arithmetic Obstructions},
submitted to Experimental Mathematics,
manuscript~ID~268704746, 2026.

\bibitem[Raayoni et~al.(2021)]{raayoni2021}
G.~Raayoni, S.~Gottlieb, Y.~Manor, G.~Pisha, Y.~Harris,
U.~Mendlovic, D.~Haviv, Y.~Hadad, and I.~Kaminer,
\textit{Generating conjectures on fundamental constants with the
Ramanujan Machine},
Nature \textbf{590} (2021), 67--73.

\bibitem[Wall(1948)]{wall1948}
H.\,S.~Wall,
\emph{Analytic Theory of Continued Fractions},
Van Nostrand, 1948; reprinted Chelsea, 1967.

\bibitem[Papanokechi(2026a)]{papanokechi2026aeal}
Papanokechi,
\textit{AEAL: Agent Epistemic Accountability Layer},
Zenodo, 2026.
\textsc{doi}: \href{https://doi.org/10.5281/zenodo.19565086}{10.5281/zenodo.19565086}.

\bibitem[Papanokechi(2026b)]{papanokechi2026ztek}
Papanokechi,
\textit{ZTEK: A Zero Trust Epistemic Kernel for AI-Assisted Research
Workflows},
Zenodo, 2026.
\textsc{doi}: \href{https://doi.org/10.5281/zenodo.19591564}{10.5281/zenodo.19591564}.

\bibitem[Papanokechi(2026c)]{papanokechi2026aibs}
Papanokechi,
\textit{AI Behavior Science --- Founding Territory Paper},
Zenodo, 2026.
\textsc{doi}: \href{https://doi.org/10.5281/zenodo.19562751}{10.5281/zenodo.19562751}.

\end{thebibliography}

\end{document}
